\documentclass{article}

\usepackage{tabularx}
\usepackage{booktabs}
\usepackage{parskip}

\setlength{\parindent}{0pt}

\title{Problem Statement and Goals\\\progname}

\author{\authname}

\date{}

\input{../Comments.text}
\input{../Common.text}

\begin{document}

\maketitle

\begin{table}[hp]
\caption{Revision History} \label{TblRevisionHistory}
\begin{tabularx}{\textwidth}{llX}
\toprule
\textbf{Date} & \textbf{Developer(s)} & \textbf{Change}\\
\midrule
2026/09/24 & Mohammad Bilal & Completed sections 1.1 \& 4 \\
\bottomrule
\end{tabularx}
\end{table}

\section{Problem Statement}

\wss{You should check your problem statement with the
\href{https://smiths.github.io/capTemplate/Checklists/ProbState-Checklist.pdf}
{problem statement checklist}}. 

\wss{You can change the section headings, as long as you include the required
information.}

\subsection{Problem}

Paediatric ophthalmology appointments can be stressful and unfamiliar for children.
Procedures such as vision testing, eye drops, and slit-lamp examinations may cause anxiety, particularly when children do not know what to expect.
This anxiety can make it difficult to cooperate during examinations, potentially affecting the quality and completeness of clinical assessments.
For example, a study of 298 children undergoing cycloplegic eye examinations found that 26\% experienced significant distress before receiving eye drops,
while uncooperative children were also more anxious and required more time to complete the examination \cite{ChiaEtAl2015}.

Children and their families therefore need adequate preparation for common ophthalmology procedures before attending an appointment.
Without appropriate preparation, unfamiliar clinical experiences may contribute to increased anxiety, reduced cooperation, and longer or more difficult appointments.
This project addresses the need for an accessible, child-friendly way to help children become more familiar and comfortable with the paediatric ophthalmology experience.

\subsection{Inputs and Outputs}

\wss{Characterize the problem in terms of ``high level'' inputs and outputs.  
Use abstraction so that you can avoid details.}

\subsection{Stakeholders}

\subsection{Environment}

\wss{Hardware and Software Environment for the users.  The developer environment
is summarized as part of the developer plan.}

\section{Goals}

\section{Stretch Goals}

\wss{Goals you hope to get to, but not necessary.  They are also helpful if the
scope of the original project needs to be expanded.}

\section{Custom Documents (CDs)}

\wss{For CAS 741: State whether the project is a research project. This
designation, with the approval (or request) of the instructor, can be modified
over the course of the term.}

\wss{For SE Capstone: List your custom documents (CDs).  Potential CDs include
usability testing, code walkthroughs, user documentation, formal proof,
GenderMag personas, Design Thinking, etc.  (The full list is on the course
outline and in Lecture 02.) Normally the number of CDs will be two (one per
term).  Approval of the CDs will be part of the discussion with the instructor
for approving the project. The CDs, with the approval (or request) of the
instructor, can be modified over the course of the term.}

\newpage{}

\section*{Appendix --- Reflection}

\wss{Not required for CAS 741}

\input{../Reflection.text}

\begin{enumerate}
    \item What went well while writing this deliverable? 
    \item What pain points did you experience during this deliverable, and how
    did you resolve them?
    \item How did you and your team adjust the scope of your goals to ensure
    they are suitable for a Capstone project (not overly ambitious but also of
    appropriate complexity for a senior design project)?
\end{enumerate}

\bibliographystyle{plain}
\bibliography{../../refs/References}

\end{document}